\documentclass[11pt]{article}
\usepackage[utf8]{inputenc}
\usepackage[margin=1in]{geometry}
\usepackage{amsmath}
\usepackage{graphicx}
\usepackage{booktabs}
\usepackage{hyperref}
\usepackage{natbib}
\usepackage{setspace}
\onehalfspacing

\title{Sex Dependence of Opioid-Mediated Responses to Subanesthetic Ketamine}
\author{Author A$^{1}$, Author B$^{1,2}$, Author C$^{2}$ \\
\small $^{1}$Department of Psychiatry, Medical School \\
\small $^{2}$Department of Pharmacology, Medical School}
\date{}

\begin{document}
\maketitle

\begin{abstract}
Ketamine produces rapid antidepressant effects, but the role of the opioid system in mediating these effects remains controversial, with conflicting clinical and preclinical findings. Here we used pharmacological functional ultrasound imaging (fUSI) in awake rats to demonstrate that opioid receptor blockade with naltrexone (10 mg/kg) suppresses ketamine-evoked cerebral blood volume (CBV) changes in a sex-dependent manner. In males, naltrexone significantly reduced ketamine-evoked CBV increases in the cingulate cortex (Cg1; Hedge's $g = -1.20$), nucleus accumbens core (NAcC; $g = -0.97$), and reversed the CBV decrease in the lateral habenula (LHb; $g = 0.87$). In females, naltrexone had no significant effect on any region. Surgical gonadectomy in males fully reversed the opioid-dependent effects, establishing that male gonadal hormones are required. Ketamine increased PSD-95 expression in male medial prefrontal cortex (blocked by naltrexone) but not in female or castrated male mPFC. Behavioral sensitization to repeated ketamine was blocked by naltrexone in males but not females. The selective NMDA receptor antagonist MK-801 showed no opioid dependence in either sex, demonstrating that the opioid component is specific to ketamine rather than a general property of NMDA receptor antagonism. These findings reveal a sex-dependent opioid mechanism in ketamine's neural and behavioral effects with implications for the therapeutic use and abuse liability of ketamine.
\end{abstract}

\section{Introduction}

Ketamine, a noncompetitive N-methyl-D-aspartate receptor (NMDAR) antagonist, has emerged as a transformative treatment for treatment-resistant depression, producing rapid antidepressant effects within hours of a single subanesthetic dose \citep{berman2000,zarate2006}. However, the mechanisms underlying ketamine's antidepressant action extend beyond simple NMDAR blockade. The AMPA receptor, BDNF-TrkB signaling, mTOR pathway, and lateral habenula burst firing have all been implicated \citep{zanos2018,autry2011,li2010,yang2018}. Among the most contentious questions is whether the opioid system plays a necessary role in ketamine's therapeutic effects.

Clinical evidence on this question is contradictory. One study reported that pretreatment with naltrexone, a non-selective opioid receptor antagonist, attenuated ketamine's antidepressant effects in treatment-resistant depression patients \citep{williams2018}. However, a subsequent randomized controlled trial found no attenuation by naltrexone \citep{yoon2019}. Preclinical evidence is similarly conflicting, with some studies showing opioid dependence of ketamine's antidepressant-like effects and others showing opioid independence \citep{zanos2019}.

A critical variable that has been minimally assessed is biological sex. Sex differences in opioid pharmacology are well established: females show greater analgesic responses to kappa-opioid agonists, different mu-opioid receptor densities, and sex-specific interactions between gonadal hormones and opioid signaling \citep{mogil2012,craft2004}. Given the extensive sex differences in both depression prevalence and opioid pharmacology, the opioid dependence of ketamine's effects may differ fundamentally between males and females.

Here we used pharmacological functional ultrasound imaging (fUSI) in awake Sprague-Dawley rats to measure the effects of opioid blockade on ketamine-evoked neural activity changes with high spatial and temporal resolution. fUSI measures cerebral blood volume (CBV) as a proxy for neural activity and provides superior sensitivity to fMRI in rodents \citep{mace2011,deffieux2018}. We combined fUSI with immunohistochemistry for structural plasticity markers, behavioral sensitization assays, and surgical gonadectomy to comprehensively characterize the sex dependence of opioid-mediated ketamine responses.

\section{Methods}

\subsection{Animals and Experimental Design}

Adult male and female Sprague-Dawley rats underwent chronic craniotomy with polymeric prosthesis for repeated fUSI. The within-subject design included three imaging sessions per animal: vehicle + ketamine (VEH+KET), naltrexone + ketamine (NTX+KET), and naltrexone + vehicle (NTX+VEH), with sessions separated by at least 48 hours.

\begin{table}[h]
\centering
\caption{Experimental design parameters.}
\label{tab:design}
\begin{tabular}{ll}
\toprule
Parameter & Value \\
\midrule
Species & Sprague-Dawley rats \\
Imaging modality & fUSI (cerebral blood volume) \\
Ketamine dose & 10 mg/kg IV (fUSI) or IP (behavior) \\
Naltrexone dose & 10 mg/kg SC \\
NTX pretreatment timing & 10 min before ketamine \\
Imaging planes & Bregma +2.5 mm, $-3.5$ mm \\
Within-subject sessions & VEH+KET, NTX+KET, NTX+VEH \\
\bottomrule
\end{tabular}
\end{table}

\subsection{Functional Ultrasound Imaging}

fUSI was performed in awake, head-fixed rats using a linear ultrasound probe positioned over the cranial window. Power Doppler images were acquired at coronal planes at bregma +2.5~mm (targeting Cg1, PrL, CPu, NAcC, NAcS) and $-3.5$~mm (targeting LHb). Baseline CBV was recorded for 10 minutes before drug injection, followed by 30 minutes of post-injection recording.

\subsection{Regions of Interest}

Seven ROIs were defined based on anatomical landmarks (Table~\ref{tab:rois}).

\begin{table}[h]
\centering
\caption{Regions of interest analyzed.}
\label{tab:rois}
\begin{tabular}{lll}
\toprule
ROI & Brain Region & Bregma Plane \\
\midrule
Cg1 & Cingulate cortex area 1 & +2.5 mm \\
PrL & Prelimbic cortex & +2.5 mm \\
CPu & Caudate-putamen & +2.5 mm \\
NAcC & Nucleus accumbens core & +2.5 mm \\
NAcS & Nucleus accumbens shell & +2.5 mm \\
LHb & Lateral habenula & $-3.5$ mm \\
LPRL & Lateral prelimbic & +2.5 mm \\
\bottomrule
\end{tabular}
\end{table}

\subsection{Gonadectomy}

A subset of male rats underwent bilateral orchiectomy under isoflurane anesthesia at least 3 weeks before fUSI experiments, allowing gonadal hormone levels to decline.

\subsection{PSD-95 Immunohistochemistry}

Separate cohorts received NTX or vehicle followed by ketamine or vehicle (10 mg/kg IP) and were perfused 24 hours later. Brain sections through mPFC were stained for PSD-95 and quantified by fluorescence intensity in defined ROIs.

\subsection{Behavioral Sensitization}

Rats received 5 consecutive daily IP ketamine injections (10 mg/kg) with or without NTX pretreatment. Locomotor activity was measured in an open field for 60 minutes following each injection. Sensitization was defined as a significant increase in locomotion from day 1 to day 5.

\subsection{Statistical Analysis}

fUSI data were analyzed using two-tailed paired $t$-tests (NTX+KET vs VEH+KET within sex) and unpaired $t$-tests (between-sex comparisons). Effect sizes were computed as Hedge's $g$. PSD-95 data were analyzed with two-way ANOVA (treatment $\times$ sex/gonadectomy). Behavioral data used repeated-measures ANOVA with Bonferroni correction.

\section{Results}

\subsection{Ketamine Produces Widespread CBV Changes in Males}

In male rats, ketamine (VEH+KET condition) produced robust CBV changes across all measured regions. Strong CBV increases were observed in Cg1 and PrL (medial prefrontal cortex), with moderate increases in NAcC, NAcS, and CPu. In contrast, the LHb showed a significant CBV decrease, consistent with ketamine's known suppressive effect on habenula activity.

\subsection{Naltrexone Suppresses Ketamine Effects in Males but Not Females}

The central finding of this study is the sex-dependent effect of opioid blockade on ketamine-evoked CBV changes (Table~\ref{tab:sexeffects}). In males, pretreatment with naltrexone significantly attenuated ketamine-evoked CBV increases in Cg1 (Hedge's $g = -1.20$, $p < 0.05$), CPu ($g = -0.78$, $p < 0.05$), NAcC ($g = -0.97$, $p < 0.05$), and LPRL ($g = 0.97$, $p < 0.05$). Naltrexone also reversed the ketamine-evoked CBV decrease in LHb ($g = 0.87$, $p < 0.05$), restoring activity toward baseline.

In females, naltrexone had no significant effect on ketamine-evoked CBV changes in any region. Effect sizes were small and non-significant (Cg1: $g = 0.53$, NS; NAcC: $g = 0.08$, NS; LHb: $g = 0.17$, NS).

\begin{table}[h]
\centering
\caption{Effect of naltrexone on ketamine-evoked CBV changes by sex (Hedge's $g$, NTX+KET vs VEH+KET).}
\label{tab:sexeffects}
\begin{tabular}{lcccc}
\toprule
ROI & Male $g$ & Male $p$ & Female $g$ & Female $p$ \\
\midrule
Cg1 & $-1.20$ & $< 0.05$ & 0.53 & NS \\
CPu & $-0.78$ & $< 0.05$ & 0.40 & NS \\
NAcC & $-0.97$ & $< 0.05$ & 0.08 & NS \\
LHb & 0.87 & $< 0.05$ & 0.17 & NS \\
LPRL & 0.97 & $< 0.05$ & $-0.05$ & NS \\
\bottomrule
\end{tabular}
\end{table}

\subsection{Gonadectomy Reverses Opioid Dependence in Males}

To determine whether the male-specific opioid dependence requires gonadal hormones, we repeated the fUSI experiments in gonadectomized males. Gonadectomy fully reversed the opioid-dependent effects: naltrexone no longer suppressed ketamine-evoked CBV changes in any region of castrated males. This demonstrates that the opioid component of ketamine's neural effects in males is dependent on male gonadal hormones, most likely testosterone or its metabolites.

\subsection{MK-801 Shows No Opioid Dependence}

To test whether the opioid dependence is a general property of NMDAR antagonism, we examined the effects of naltrexone on MK-801 (a more selective NMDAR antagonist). Naltrexone did not suppress MK-801-evoked CBV changes in either sex, demonstrating that the opioid-dependent component is specific to ketamine and not a consequence of NMDAR blockade per se.

\subsection{PSD-95 Structural Plasticity Shows Same Sex-Opioid Pattern}

Ketamine increased PSD-95 immunoreactivity in mPFC of intact males 24 hours after administration, consistent with enhanced synaptogenesis. This increase was fully blocked by naltrexone pretreatment. In females, ketamine did not significantly increase PSD-95, and in castrated males, ketamine also failed to increase PSD-95. The pattern of PSD-95 changes precisely mirrored the fUSI findings, suggesting a common opioid-dependent mechanism underlying both acute neural activity changes and structural plasticity.

\subsection{Behavioral Sensitization Is Opioid-Dependent in Males Only}

Repeated ketamine administration (5 daily injections) produced locomotor sensitization in both males and females. However, naltrexone pretreatment blocked sensitization only in males; female sensitization was unaffected by naltrexone (Table~\ref{tab:behavior}).

\begin{table}[h]
\centering
\caption{Behavioral sensitization to repeated ketamine by sex and naltrexone treatment.}
\label{tab:behavior}
\begin{tabular}{llcc}
\toprule
Group & Treatment & Sensitization & NTX Blocked? \\
\midrule
Male & VEH+KET & Present & --- \\
Male & NTX+KET & Absent & Yes \\
Female & VEH+KET & Present & --- \\
Female & NTX+KET & Still present & No \\
\bottomrule
\end{tabular}
\end{table}

\section{Discussion}

Our findings demonstrate that the opioid system plays a sexually dimorphic role in ketamine's neural, structural, and behavioral effects. In males, opioid receptor blockade suppresses ketamine-evoked activity changes in depression-relevant brain regions, blocks ketamine-induced synaptogenesis, and prevents behavioral sensitization. In females and castrated males, these opioid-dependent effects are absent, indicating a requirement for male gonadal hormones.

The sex dependence of opioid involvement may help reconcile the contradictory clinical evidence regarding naltrexone's effect on ketamine's antidepressant action. The clinical trial reporting attenuation by naltrexone enrolled predominantly male participants \citep{williams2018}, while studies finding no attenuation may have had more balanced sex ratios. Our results predict that naltrexone would attenuate ketamine's antidepressant effects preferentially in males, a prediction that could be tested in sex-stratified clinical trials.

The gonadectomy reversal demonstrates that the opioid-dependent component requires circulating gonadal hormones rather than being determined by chromosomal sex. Testosterone modulates opioid receptor expression and function through multiple mechanisms, including direct transcriptional regulation and indirect effects via aromatization to estradiol \citep{craft2004}. The specific hormonal mechanism connecting gonadal hormones to the opioid-ketamine interaction remains to be elucidated.

The dissociation between ketamine and MK-801 is particularly informative. MK-801 is a more selective and potent NMDAR antagonist than ketamine, yet its effects show no opioid dependence. This suggests that the opioid component of ketamine's action arises from a mechanism distinct from NMDAR blockade, potentially involving direct or indirect modulation of opioid receptors or opioid peptide release \citep{gupta2011,bhatt2020}.

The convergence of fUSI, PSD-95, and behavioral findings strengthens the conclusion that opioid dependence is a fundamental feature of ketamine's mechanism in males. The same sex-specific opioid pattern was observed across three independent measures spanning acute neural activity (fUSI), structural plasticity (PSD-95), and behavior (locomotor sensitization). This convergence argues against confounds specific to any single technique.

These findings have important implications for ketamine's abuse liability. Opioid-dependent rewarding effects in males but not females suggest that addiction risk may be sex-dependent, a consideration relevant to the expanding clinical use of ketamine and esketamine during the ongoing opioid crisis \citep{schak2016}.

\section{Conclusion}

We have demonstrated that opioid receptor blockade with naltrexone suppresses ketamine-evoked neural activity changes, structural plasticity, and behavioral sensitization in a sex-dependent manner, with effects present in males but absent in females and castrated males. These findings establish a gonadal hormone-dependent opioid component in ketamine's mechanism of action and provide a framework for reconciling contradictory clinical evidence regarding the role of opioids in ketamine's antidepressant effects.

\bibliographystyle{plainnat}
\bibliography{references}

\end{document}
